% GENERATED FILE. Edit canonical lessons, then run npm run generate:books.
% canonical-source-hash: fnv1a64:64631c0852b42e9b
% canonical-lessons: LA-C09-anni-tempora

\chapter{Seasons}
\label{ch:la-seasons}
\hlchaptermodality{9}

\begin{chapteropening}
\textbf{By the end of this chapter:} \emph{I can name the four seasons in Latin and say which ones Romance rebuilt out of a phrase instead.}

Name vēr, aestās, autumnus, and hiems, then sort which survived plainly and which Romance rebuilt from phrases: \textquotedblleft{}first spring\textquotedblright{} and hībernum.
\end{chapteropening}

\section[vēr, aestās, autumnus, hiems]{vēr, aestās, autumnus, hiems — Latin's own four seasons}

\label{lesson:LA-C09-anni-tempora}



\begin{quote}
Four Latin nouns, one for each season — and a genuine surprise about which ones actually survived into Spanish and French as ordinary words.
\end{quote}

\subsection*{You'll want to know: The four seasons, plain and simple}
\noindent
\begin{tabularx}{\linewidth}{@{}>{\raggedright\arraybackslash}X>{\raggedright\arraybackslash}X@{}}
\toprule
\textbf{Latin} & \textbf{meaning} \\
\midrule
\textbf{vēr, vēris} & \textbf{spring} \\
\textbf{aestās, aestātis} & \textbf{summer} (also just \textquotedblleft{}heat\textquotedblright{}) \\
\textbf{autumnus, autumnī} & \textbf{autumn} \\
\textbf{hiems, hiemis} & \textbf{winter} \\
\bottomrule
\end{tabularx}

\begin{culture}
\textbf{autumnus} survived plainly: it's the direct ancestor of Spanish \emph{otoño}, French \emph{automne}, and English \emph{autumn} itself — no detours.

But \textbf{vēr} and \textbf{hiems}, the plain nouns for spring and winter, mostly \textbf{didn't} make it into everyday Romance vocabulary. Instead, Spanish and French built brand-new words out of \textbf{adjective phrases} describing those seasons.

Spring was not \emph{vēr} alone, but a Vulgar Latin phrase/reanalysis, roughly \textquotedblleft{}\textbf{first spring}\textquotedblright{}:

\textbf{prīma} (\textquotedblleft{}first\textquotedblright{}) + \textbf{vēra} (\emph{vēr}, reshaped) $\to$ \textbf{primavera}

That gives Spanish \emph{primavera}, Italian \emph{primavera} (still the everyday word there), and French (originally) \emph{primevère} — later displaced there by \emph{printemps} (\textquotedblleft{}first time,\textquotedblright{} the same \textquotedblleft{}first\textquotedblright{} put with \emph{tempus}: the very same \textquotedblleft{}adjective + tempus\textquotedblright{} recipe as winter's word below). \emph{Primevère} narrowed to just mean \textquotedblleft{}\textbf{primrose}\textquotedblright{} today.

Winter was not \emph{hiems} alone, but \textbf{hībernum (tempus)}, \textquotedblleft{}\textbf{the wintry {[}time{]}}\textquotedblright{} $\to$ Spanish \emph{invierno}, French \emph{hiver} — the same root as English \textbf{hibernate}.

\textbf{Aestās} (\textquotedblleft{}summer/heat\textquotedblright{}) did survive directly into French \emph{été} — but Spanish took yet another path: it built \textbf{vērānum (tempus)} (\textquotedblleft{}of spring\textquotedblright{}) from \emph{vēr} itself, and that word's \emph{meaning} later drifted forward to name \textbf{summer} instead of spring, giving Spanish \emph{verano} — leaving \emph{primavera} to do spring's job.

So the picture is mixed, not a clean split: \textbf{autumnus} cruised straight through into Spanish, French, \emph{and} English unchanged; \textbf{aestās} cruised through into French (\emph{été}) but got replaced in Spanish; \textbf{vēr} and \textbf{hiems}, the plain nouns for spring and winter, essentially left \textbf{no} everyday trace in either language — both were replaced by phrase-built words instead (or, in \emph{verano}'s case, by a related word whose meaning simply moved).
\end{culture}

\subsection*{Your turn}
Say these aloud:

\begin{itemize}
\raggedright
  \item \textquotedblleft{}vēr, aestās, autumnus, hiems\textquotedblright{} — the four seasons
  \item the one that survived plainly — \textquotedblleft{}autumnus $\to$ otoño, automne, autumn\textquotedblright{}
  \item the two built from phrases — \textquotedblleft{}first spring\textquotedblright{} $\to$ primavera, \textquotedblleft{}hībernum $\to$ invierno, hiver\textquotedblright{}
\end{itemize}

\begin{tcolorbox}[breakable,colback=teal!4,colframe=teal!35!black,title={Before you move on}]
What are Latin's four season-nouns? (\textbf{Vēr, aestās, autumnus, hiems}.) Which one survived plainly into Spanish, French, and English? (\textbf{Autumnus} $\to$ \emph{otoño, automne, autumn}.) What did Spanish and French build instead of the plain words \emph{vēr} and \emph{hiems}? (\textbf{Adjective phrases} — \textquotedblleft{}first spring\textquotedblright{} $\to$ \emph{primavera/primevère}; \emph{hībernum} \textquotedblleft{}wintry\textquotedblright{} $\to$ \emph{invierno/hiver}.)
\end{tcolorbox}
